\documentclass[11pt]{article}
\usepackage[margin=1in]{geometry}
\usepackage{amsmath,amssymb,hyperref}
\title{A counterexample to Conjecture 00000008554}
\author{gaochengzhecpu}
\date{}
\begin{document}
\maketitle

The asserted homological dimension is false already for the Boolean lattice
$B_2$.  The counterexample works with both ordinary and reduced integral
homology.  It also identifies precisely the Euler-characteristic convention:
ordinary Euler characteristic is not the M\"obius number, whereas reduced Euler
characteristic is.

\paragraph{The lattice and its crosscut.}
Let $L=\{0,a,b,1\}$, with $0<a<1$ and $0<b<1$, and with $a,b$
incomparable.  This is the lattice of subsets of a two-element set.  Its rank
is $2$: every maximal chain has three elements and two strict steps.
The set $C=\{a,b\}$ is an antichain in the proper part of $L$ and meets
every maximal chain.  It is simultaneously the set of atoms and the set of
coatoms.

Use the standard crosscut complex, whose faces are the nonspanning subsets
$S\subseteq C$: equivalently, $\bigwedge S\ne0$ or $\bigvee S\ne1$.
Here the empty face and the two singletons are faces, but $\{a,b\}$ is not,
because $a\wedge b=0$ and $a\vee b=1$.  Thus the crosscut complex $K$ is
exactly two isolated vertices.  The atom and coatom crosscut conventions give
the same complex in this example.  For the standard definitions and the
reduced Euler convention, see Section 2 of T. McConville,
\emph{Crosscut-simplicial lattices},
\href{https://arxiv.org/abs/1409.6269}{arXiv:1409.6269}.

\paragraph{Homology in all degrees.}
The integral simplicial chain groups and boundaries are
\[
 C_0(K;\mathbb Z)=\mathbb Z a\oplus\mathbb Z b,\qquad
 C_n(K;\mathbb Z)=0\quad(n\ge1).
\]
Consequently $H_0(K;\mathbb Z)\cong\mathbb Z^2$ and
$H_n(K;\mathbb Z)=0$ for every $n\ge1$.
For reduced homology, the augmentation sends $u a+v b$ to $u+v$.
Its kernel is generated by $a-b$, so
$\widetilde H_0(K;\mathbb Z)\cong\mathbb Z$ and
$\widetilde H_n(K;\mathbb Z)=0$ for every $n\ge1$.
In either convention, the largest degree of nonzero homology is therefore
$0$, while the conjectured value is
\[
 \operatorname{rank}(L)-1=2-1=1.
\]
This disproves the dimension assertion and hence the stated conjunction.
The simplicial dimension of $K$ is also $0$; it cannot rescue the asserted
value.  No claim about any different algebraic invariant called a dimension
is needed.

\paragraph{Euler characteristic.}
The M\"obius recurrence gives
$\mu(0,0)=1$, $\mu(0,a)=\mu(0,b)=-1$, and $\mu(0,1)=1$.
The complex has two vertices and no higher nonempty faces.  Thus
$\chi(K)=2\ne1=\mu(0,1)$, but
$\widetilde\chi(K)=\chi(K)-1=1=\mu(0,1)$.
Accordingly the Euler clause fails with ordinary Euler characteristic, and
holds in this example with reduced Euler characteristic.  The dimension
counterexample is independent of that convention.

\paragraph{Formal verification and its scope.}
The accompanying Lean 4.19.0 file imports only \texttt{Std}.  It checks the
lattice laws, maximum and maximal chain lengths, and the crosscut property.
Every actual predicate subset of the lattice is represented by a finite
code; inclusion-maximal chains are transferred back to arbitrary predicate
subsets.  Every subset of the crosscut is likewise represented.  Face meets
and joins are verified by their universal properties, so the face list is
exhaustive.

For every natural degree $n$, the file defines the full type of oriented
$n$-simplices and its group of integer coefficient functions.  It proves
there are no simplices in any positive degree, hence every positive chain
group is zero.  The theorem \texttt{boundary\_forced} shows that every
zero-preserving candidate for the simplicial boundary must be the displayed
zero map.  Nonzero homology is expressed as a cycle that is not a boundary,
and zero homology as every cycle being a boundary; these are precisely the
nonzero and zero conditions for the quotient of cycles by boundaries.
The witnesses $(1,0)$ and $(1,-1)$ establish the ordinary and reduced claims,
respectively.  There is no bounded-degree search in the homology argument.
The main theorem is
\begin{center}
\texttt{Conjecture8554.conjecture8554\_counterexample}.
\end{center}
Separate theorems verify the M\"obius recurrence, its uniqueness for arbitrary
integer-valued functions, and both Euler conventions.  The file uses no
\texttt{sorry}, \texttt{admit}, \texttt{native\_decide}, or custom axioms;
its final commands report the axioms used.
\end{document}
